%% Requires compilation with XeLaTeX or LuaLaTeX
\documentclass[10pt,xcolor={table,dvipsnames},t]{beamer}
\usetheme{UVic}

\usepackage{booktabs}

\title[PhD Progress Report]{PhD Progress Report}
\subtitle{Learning Automata Framework for Adaptive Kelp Forest Monitoring in British Columbia Coastal Waters}
\author{Ekaba Bisong}
\institute{Department of Computer Science\\University of Victoria}
\date{January 12, 2025}

\begin{document}

% Slide 1: Title
\begin{frame}
  \titlepage
\end{frame}

% Slide 2: Progress Overview
\begin{frame}{Progress Overview}

\textbf{Since December 2025:}
\begin{itemize}
  % \item Thesis proposal approved by committee
  \item Began implementation of Paper 1: ``Learning Automata for Adaptive Kelp Detection''
  \item Created publication-quality maps for study area visualization
  \item Established code infrastructure for visualization and analysis
\end{itemize}

\vspace{0.5em}
\textbf{Current Focus:}
\begin{itemize}
  \item Paper 1 methodology implementation
  \item Study area maps and figures for publication
  \item Data preprocessing pipeline refinement
\end{itemize}

\end{frame}

% Slide 3: Paper 1 Status
\begin{frame}{Paper 1: Current Status}

\textbf{Title:} ``Learning Automata for Adaptive Kelp Detection with Limited Labels''

\vspace{0.5em}
\textbf{Completed:}
\begin{itemize}
  \item Study area map (BC coastal waters with Sentinel-2 tile footprints)
  \item Data distribution map (kelp coverage by scene)
  \item Environmental context map (bathymetry, substrate, temperature gradient)
  \item Train/test split visualization (Protocol B: spatial adaptation)
  \item Sentinel-2 mosaic with RGB sample composites
\end{itemize}

\vspace{0.5em}
\textbf{In Progress:}
\begin{itemize}
  \item Baseline heuristic implementations (spectral indices)
  \item LA framework core components
\end{itemize}

\end{frame}

% Slide 4: Study Area Coverage
\begin{frame}{Study Area: BC Coastal Waters}

\textbf{Sentinel-2 Tile Coverage:}
\begin{itemize}
  \item 7 scenes spanning BC's coastline (48.7°N -- 53.0°N)
  \item 416 tiles total (512$\times$512 pixels, 10m resolution)
  \item Temporal span: 2021--2023
\end{itemize}

\vspace{0.5em}
\begin{table}[h]
\small
\begin{tabular}{lll}
\toprule
\textbf{Location} & \textbf{Year} & \textbf{Tiles} \\
\midrule
S. Vancouver Island & 2021 & 42 \\
Salish Sea & 2022 & 72 \\
Strait of Georgia & 2022 & 73 \\
C. Vancouver Island & 2023 & 58 \\
Discovery Islands & 2023 & 56 \\
Central Coast & 2023 & 59 \\
North Coast & 2023 & 56 \\
\bottomrule
\end{tabular}
\end{table}

\end{frame}

% Slide 5: Data Statistics
\begin{frame}{Data Distribution Statistics}

\textbf{Kelp Coverage Analysis:}
\begin{itemize}
  \item Overall kelp coverage: $\sim$0.84\% of pixels
  \item Range across scenes: 0.20\% -- 1.68\%
  \item Highly imbalanced class distribution (typical for kelp detection)
\end{itemize}

\vspace{0.5em}
\textbf{Environmental Heterogeneity:}
\begin{itemize}
  \item Bathymetry variation: 10.3 -- 218.4 (normalized values)
  \item Substrate variation: 1.1 -- 2.7 (categorical encoding)
  \item Latitudinal gradient: 4.3° span (temperature proxy)
\end{itemize}

\vspace{0.5em}
\textbf{Implication:} Environmental variability motivates context-aware adaptation in the LA framework.

\end{frame}

% Slide 6: Visualization Infrastructure
\begin{frame}{Visualization Infrastructure}

\textbf{Publication-Quality Maps Created:}
\begin{enumerate}
  \item \texttt{fig1\_study\_area} --- Main study area with tile locations
  \item \texttt{fig2\_data\_distribution} --- Kelp coverage and tile counts
  \item \texttt{fig3\_environmental\_context} --- Bathymetry, substrate, zones
  \item \texttt{fig4\_sentinel\_mosaic} --- RGB composites with samples
  \item \texttt{fig\_train\_test\_split} --- Protocol B spatial split
\end{enumerate}

\vspace{0.5em}
\textbf{Technical Stack:}
\begin{itemize}
  \item Python with Cartopy, Matplotlib, GeoPandas, Rasterio
  \item BC Albers projection (EPSG:3005)
  \item 300 DPI output for journal submission
\end{itemize}

\end{frame}

% Slide 7: Protocol B: Spatial Adaptation
\begin{frame}{Protocol B: Spatial Adaptation Experiment}

\textbf{Train/Test Split Design:}
\begin{itemize}
  \item \textbf{Training (2021--2022):} 3 scenes, 187 tiles
  \begin{itemize}
    \item S. Vancouver Island, Salish Sea, Strait of Georgia
  \end{itemize}
  \item \textbf{Test (2023):} 4 scenes, 229 tiles
  \begin{itemize}
    \item C. Vancouver Island, Discovery Islands, Central Coast, North Coast
  \end{itemize}
\end{itemize}

\vspace{0.5em}
\textbf{Research Question Addressed (RQ1):}
\begin{quote}
How effectively can an LA framework adapt to new geographic regions with limited labeled examples?
\end{quote}

\vspace{0.3em}
\textbf{Evaluation:} Zero-shot and few-shot adaptation to unseen 2023 regions.

\end{frame}

% Slide 8: Implementation Roadmap
\begin{frame}{Implementation Roadmap}

\textbf{Phase 1: Foundation (Current)}
\begin{itemize}
  \item[$\checkmark$] Data preprocessing and tile generation
  \item[$\checkmark$] Study area visualizations
  \item[$\square$] Tier 1 spectral indices (NDVI, FAI, NDWI, Kelp Index)
  \item[$\square$] Baseline evaluation framework
\end{itemize}

\vspace{0.5em}
\textbf{Phase 2: Core Framework (Next)}
\begin{itemize}
  \item[$\square$] Component 1: Global Policy Learning (L$_{\text{R-I}}$)
  \item[$\square$] Component 2: Context-Aware Local Adaptation
  \item[$\square$] Component 3: Sparse Label Handling
\end{itemize}

\vspace{0.5em}
\textbf{Phase 3: Experiments}
\begin{itemize}
  \item[$\square$] Protocol A: Temporal generalization
  \item[$\square$] Protocol B: Spatial adaptation
  \item[$\square$] Protocol C: Sparse labels
\end{itemize}

\end{frame}

% Slide 9: Next Steps
\begin{frame}{Immediate Next Steps}

\textbf{This Week (Jan 13--19):}
\begin{itemize}
  \item Implement Tier 1 spectral indices as baseline heuristics
  \item Set up evaluation metrics (IoU, F1, Precision, Recall)
  \item Begin LA global policy component
\end{itemize}

\vspace{0.5em}
\textbf{This Month (January 2025):}
\begin{itemize}
  \item Complete baseline experiments on training set
  \item Implement context feature extraction (environmental vectors)
  \item Draft Paper 1 introduction and related work sections
\end{itemize}

\vspace{0.5em}
\textbf{Target:} Paper 1 submission by November 2026 (per timeline)

\end{frame}

% Slide 10: Summary
\begin{frame}{Summary}

\begin{itemize}
  \item \textbf{Proposal:} Approved --- now in implementation phase

  \vspace{0.5em}
  \item \textbf{Paper 1 Progress:} Study area visualizations complete; baseline implementation underway

  \vspace{0.5em}
  \item \textbf{Data:} 416 tiles across 7 scenes ready for experiments

  \vspace{0.5em}
  \item \textbf{Infrastructure:} Visualization pipeline established with publication-quality outputs

  \vspace{0.5em}
  \item \textbf{Next Milestone:} Baseline heuristics and evaluation framework
\end{itemize}

\vspace{1em}
\centering
\textbf{Questions?}

\end{frame}

\end{document}
